\documentclass[12pt,letterpaper, onecolumn]{exam}
\usepackage{amsmath}
\usepackage{amssymb}
\usepackage{amsthm}
\usepackage{graphicx}
\usepackage{mathrsfs}
\usepackage{xcolor}
\usepackage{ulem}
\usepackage[lmargin=71pt, tmargin=1.2in]{geometry}  %For centering solution box
\usepackage[english]{babel}
\newtheorem{theorem}{Theorem}
\newtheorem{lemma}[theorem]{Lemma}
\usepackage{accents}
\newcommand{\ubar}[1]{\underaccent{\bar}{#1}}
\newcommand{\bs}[1]{\boldsymbol{#1}}
\newcommand{\bm}[1]{\mathbf{#1}}
\newcommand{\der}[2]{\frac{d #1}{d#2}}
\newcommand{\pder}[3]{\frac{\delta^{#3} #1}{\delta^{#3}#2}}
\newcommand{\rs}{$X_1, ..., X_n$}
\newcommand{\nsum}{\sum_{i=1}^n}
\newcommand{\nprod}{\prod_{i=1}^n}
\newcommand*{\Comb}[2]{{}^{#1}C_{#2}}%
\newcommand{\ifc}{\textrm{ if }}
\newcommand{\real}{\mathbb{R}}
\newcommand{\nat}{\mathbb{N}}
\newcommand{\power}{\mathcal{P}}
\newcommand{\rational}{\mathbb{Q}}
\newcommand{\cutelim}{\lim\limits_{n\to\infty}}
\newcommand{\ibold}{\mathbf{I}}
\newcommand{\Sbb}{\mathbb{S}}
\newcommand{\sigmaF}{\mathcal{F}}
\newcommand{\sigmaG}{\mathcal{G}}
\newcommand{\calA}{\mathcal{A}}
\newcommand{\calO}{\mathcal{O}}
\newcommand{\calC}{\mathcal{C}}
\newcommand{\calB}{\mathcal{B}}
\newcommand{\calI}{\mathcal{I}}
\newcommand{\salg}{$\sigma$-algebra }
\newcommand{\gensig}[1]{\sigma\langle #1 \rangle} % notation for a sigma algebra generated by some set 
\newcommand{\om}{\mu^*} % outer measure 
\newcommand{\rat}{\mathbb{Q}}


% \chead{\hline} % Un-comment to draw line below header
\thispagestyle{empty}   %For removing header/footer from page 1

\begin{document}

\begingroup  
    \centering
    \LARGE STAT 6410\\
    \LARGE Homework 3 \\[0.5em]
    \large Sarah Kilgard\par
\endgroup
\rule{\textwidth}{0.4pt}
\pointsdroppedatright   %Self-explanatory
\printanswers
\renewcommand{\solutiontitle}{\noindent\textbf{Ans:}\enspace}   %Replace "Ans:" with starting keyword in solution box

\begin{questions}
%%%%%%%%%%% Question 1 %%%%%%%%%%%%%%%%%%%%    
    \question Chapter 1, Problem 1.3
    \begin{itemize}
        \item[(i)] Show that $\sigmaF_6 = \{A \in \Omega: A \textrm{ is countable or } A^c \textrm{ is countable.}\}$ from example 1.1.4 is a $\sigma$-algebra 
        \item[(ii)] Define $\mu$ on $\sigmaF_6$ as $\mu(A) = 1$ if $A$ is uncountable and $\mu(A) = 0$ if $A$ countable. Is $\mu$ a measure on $\sigmaF_6$ ? Justify.
    \end{itemize}
    \begin{solution}
    
    (i)
    \begin{proof}
    Let $\Omega$ be a non-empty set. 
    \begin{itemize}
        \item[(a)] $\Omega \subset \Omega$, and $\Omega^c = \emptyset$, which is countable (even if $\Omega$ is not). Thus, $\Omega \in \sigmaF_6$.
        \item[(b)] Let $A \in \sigmaF_6$ be given. \\
        \textit{Case 1:} $A$ is countable. \\ 
        $A^c \in \sigmaF_6$ since $(A^c)^c = A$ is countable. 
        
        \textit{Case 2:} $A^c$ is countable. \\
        $A^c \in \sigmaF_6$ by the definition of $\sigmaF_6$.
        \item[(c)] Let $A, B \in \sigmaF_6$. \\
        \textit{Case 1:} $A$ and $B$ are countable. \\ $A \cup B$ is clearly countable. Thus, $A \cup B \in \sigmaF_6$.

        \textit{Case 2:} Without loss of generality, $A^c$ and $B$ are countable. \\
        We don't know for sure if $A^c \cup B$ is countable since we don't know if $A$ is countable. However, $(A \cup B)^c = A^c \cap B^c$ is countable since we know that $A^c$ is countable. Thus, $A \cup B \in \sigmaF_6$.

        \textit{Case 3:} $A^c$ and $B^c$ are countable. \\
        We don't know for sure if $A^c \cup B$ is countable since we don't know if $A$ is countable (or $B$ is countable). However, $(A \cup B)^c = A^c \cap B^c$ is countable since $A^c$ and $B^c$ are countable. Thus, $A \cup B \in \sigmaF_6$. 
        
     Thus, $\sigmaF_6$ is an algebra.

        \item[(d)] Let $n \in \nat$ be given. \\
        Suppose $A_1, A_2, ... , A_n \in \sigmaF_6$.

        \textit{Case 1:} $A_i$ is countable for all $1 \leq i \leq n$. \\
        It is clear to see that $\bigcup_{i =1}^n A_i$ is countable. Thus, $\bigcup_{i =1}^n A_i \in \sigmaF_6$.

        \textit{Case 2:} There exists at least one $1 \leq j \leq n$ such that $A_j^c$ is countable. \\
        Since $A_j^c$ is countable, we don't know if $A_j$ is countable. Thus, we don't know if $\bigcup_{i =1}^n A_i$ is countable. \\
        However, $\left(\bigcup_{i =1}^n A_i\right)^c = \bigcap_{i =1}^n A_i^c = \bigcap_{i =1}^{j-1} A_i^c \bigcap A_j^c \bigcap_{k = j+1}^{n} A_k^c$. \\
        Since $A_j^c$ is countable, $\left(\bigcup_{i =1}^n A_i\right)^c$ is countble, and thus, $\left(\bigcup_{i =1}^n A_i\right)^c \in \sigmaF_6$
    \end{itemize}
    Thus, $\sigmaF_6$ is a $\sigma$-algebra.
    \end{proof}
    (ii) 

    \begin{enumerate}
        \item[(a)] WTS $\mu(A) \in [0, \infty] \: \forall \: A \in \sigmaF_6$

        Since $\mu(A) \in \{0, 1\} \subset [0, \infty]$, (a) holds
        \item[(b)] WTS: $\mu(\emptyset) = 0$

        $\emptyset$ is countable, and thus $\mu(\emptyset) = 0$, so (b) holds. 
        
        \item[(c)] For any disjoint collection of sets $A_1, A_2, ... \in \sigmaF_6$ with $\cup_{n \geq 1} A_n \in \sigmaF_6$, $\mu \left(\bigcup_{n \geq 1} A_n\right) = \sum_{n=1}^\infty \mu(A_n)$
        
        (c) does not hold!

        If $A_1, A_2 \in \sigmaF_6$ and both are uncountable, $\cup_{i \in \{1, 2\}} A_i$ uncountable. So, $\mu\left(\cup_{i \in \{1, 2\}} A_i\right) = 0$, however $\sum_{i=1}^2 \mu(A_i) = 1 + 1 = 2$. $1 \neq 2$, so (c) does NOT hold. 

        Thus, $\mu$ is not a measure on $\sigmaF$.
    \end{enumerate}
    \end{solution}

%%%%%%%%%%% Question 2 %%%%%%%%%%%%%%%%%%%%    
    \question Chapter 1, Problem 1.8 \\
    Let $\Omega \equiv \{1, 2,...\} = \nat$ and $A_i \equiv\{j : j \in \nat,j \geq i\}$, $i \in \nat$. Show that $\sigma \langle \calA\rangle  = \power(\Omega)$ where $\calA = \{A_i : i \in \nat\}$. 
    \begin{solution}
    \begin{proof}
        \textbf{Part 1:} WTS: $\gensig{\calA} \subset \power(\Omega)$ \\
        Since $A_i \in \Omega$ for all $i \in \nat$ which implies $\calA \subset \Omega$, it is clear to see that $\gensig{\calA} \subset \power(\Omega)$

        \textbf{Part 2:} WTS: $\gensig{\calA} \supset \power(\Omega)$

        Let $B \in \power(\Omega)$ be given. \\
        Thus, there exists some $\mathcal{J} \in \nat$ such that $B = \bigcup_{j \in \mathcal{J}} A_j$.

        Since $\gensig{\calA}$ (a) $\calA \subset \Omega$ and (b) since $\gensig{\calA}$ is a \salg, $B_n \in \gensig{\calA} \textrm{ for } n \geq 1 \implies \cup_{n \geq 1} B_n \in \gensig{\calA}$.

        Since $A_j \in \gensig{\calA}$ for all $j \in \mathcal{J}$,  $B = \bigcup_{j \in \mathcal{J}} A_j \in \gensig{\calA}$.

        Thus, $\gensig{\calA} = \power(\Omega)$
    \end{proof}        
    \end{solution}

%%%%%%%%%%% Question 3 %%%%%%%%%%%%%%%%%%%%    
    \question Chapter 1, Problem 1.10 \\
     Show that in Example 1.1.6, $\calO_j \subset \sigma \langle \calO_i\rangle $ for all $1 \leq i$, $j \leq 4$. 

     \begin{itemize}
         \item $\calO_1 = \{(a_1, b_1) \times (a_2, b_2): -\infty \leq a_1 < b_1 \leq \infty, \: -\infty \leq a_2 < b_2 \leq \infty\}$
         \item $\calO_2 = \{(-\infty, x_1) \times (-\infty, x_2): -\infty < x_1, x_2 < \infty\}$
         \item $\calO_3 = \{(a_1, b_1) \times (a_2, b_2): a_1 < b_1 , \: a_2 < b_2, a_1, a_2, b_1, b_2, \in \rat \}$
         \item $\calO_4 = \{(-\infty, x_1) \times (-\infty, x_2): x_1, x_2 \in \rat \}$
     \end{itemize}

     (Okay to do for $k = 2$, i.e. for $\real^2$.)
    \begin{solution}
    \begin{proof}
        For $k = 2$:\\
        \textbf{Part 1:} WTS: $\calO_3 \subset \gensig{\calO_1}$ \\
        We know that $\gensig{\calO_1} = \gensig{\calO_3}$, from example 1.1.6 which implies that \\ $\calO_3 \subset \gensig{\calO_3} = \gensig{\calO_1}$. 
        
        \textbf{Part 2:} WTS: $\calO_2 \subset \gensig{\calO_1}$

        By the definition of $\calO_1$ and $\calO_2$, $\calO_2 \subset \calO_1$. Since, $\calO_1 \subset \gensig{\calO_1}$, $\calO_2 \subset \gensig{\calO_1}$.

        (This is because the lower bound of the open sets in $\calO_1$ ($a_1$ and $a_2$) can be $-\infty$). 

        \textbf{Part 3:} WTS: $\calO_4 \subset \gensig{\calO_1}$

        By the same logic as \textbf{Part 2}, $\calO_4 \subset \gensig{\calO_3}$. Since $\gensig{\calO_3} = \gensig{\calO_1}$, this means that $\calO_4 \subset \gensig{\calO_1}$.
    \end{proof}
    \end{solution}

%%%%%%%%%%% Question 4 %%%%%%%%%%%%%%%%%%%%    
    \question Chapter 1, Problem 1.18 \\
    Let $\Omega \equiv \nat$, $\sigmaF = \power(\Omega)$, and $A_n = \{j : j \in \nat,j \geq n\}$, $n \in \nat$. Let $\mu$ be the counting measure on $(\Omega, \sigmaF)$. Verify that $\cutelim \mu(A_n) \neq \mu(\cap_{n\geq 1} A_n)$.

    \begin{solution}
    First, it is important to note that for all $n \in \nat$, $\mu(A_n) = \infty$. This is clear to see if we write $A_n $ like this: $A_n = \{n, n+1, n+2, n+3, ... \}$. In other words, there are always an infinite number of natural numbers greater than any given natural number, regardless of how large that given natural number is. 

    Thus, it is clear to see that: $\cutelim \mu(A_n) = \infty$.

    $\cap_{n\geq 1} A_n = \emptyset$, and by the definition of a measure, $\mu(\cap_{n\geq 1} A_n) = \mu(\emptyset) = 0.$

    Therefore,  $\cutelim \mu(A_n) \neq \mu(\cap_{n\geq 1} A_n)$
    \end{solution}

%%%%%%%%%%% Question 5 %%%%%%%%%%%%%%%%%%%%    
    \question Chapter 1, Problem 1.19 \\
    Let $\Omega$ be a nonempty set and let $\calC \subset \power(\Omega)$ be a semialgebra. Let
    \[\calA(\calC) \equiv \left\{A: A = \bigcup_{i=1}^k B_i: B_i \in \calC, \: i = 1, 2, ... , k, \: k \in \nat \right \}\]
    \begin{itemize}
        \item[(i)] Show that $\calA(\calC)$ is the smallest algebra containing $\calC$. 
        \item[(ii)] Show also that $\sigma \langle \calC\rangle  = \sigma \langle \calA(\calC)\rangle $.
    \end{itemize}

    \begin{solution}
    
    (i.)
    \begin{proof}
    \textbf{Part 1:} WTS: $\calC \subset \calA(\calC)$\\
    It is clear to see from the definition of $\calA(\calC)$ that $A \in \calC \: \implies \: A \in \calA(\calC)$. Thus $\calC \subset \calA(\calC)$
    
    \textbf{Part 2:} First, we must show that $\calA(\calC)$ is an algebra.
    \begin{itemize}
       \item[(b)] WTS: $A \in \calA(\calC)  \: \implies A^c \in \calA(\calC)$ 

        Since $A \in \calA(\calC)$, $A \in \calC$. Thus, there exists $B_1, ... , B_k \in \calC$ such that $B_i \cap B_j = \emptyset$ and $\cup_{i=1}^k B_i = A^c$. 

        Since $B_1, ... , B_k \in \calC$, by definition, $\cup_{i = 1}^k B_i \in \calA(\calC)$.

        \item[(a)]  WTS: $\Omega \in \calA(\calC)$

        Let $A \in \calA(\calC)$. By (b), $A^c \in \calA(\calC)$. 

        By definition, $A \cup A^c = \Omega \in \calA(\calC)$. 
        \item[(c)] WTS: $A, B \in \calA(\calC) \: \implies \: A \cap B \in \calA(\calC)$ 

        Let $A, B \in \calA(\calC)$ be given. Thus, $A, B \in \calC.$ So, by definition of semialgebra, $A \cap B \in  \calC \subset  \calA(\calC)$. 
    \end{itemize}
    Thus, $\calA(\calC)$ is an algebra. 

    \textbf{Part 3:} WTS: $\calA(\calC)$ is the \textbf{smallest} algebra containing $\calC$

    Let $\mathcal{E}$ be an arbitrary algebra on $\Omega$ such that $\calC \subset \mathcal{E}$.

    Let $A \in \calA(\calC)$ be given. Thus, $A = \cup_{i = 1}^k B_i$, where $\{B_i\}_{0 \leq i \leq k} \subset \calC$. By the definition of an $\mathcal{E}$, $B_i \in \mathcal{E}$ for all $0 \leq i \leq k$. By the definition of an algebra (and since $\mathcal{E}$ is an algebra), $A = \cup_{i = 1}^k B_i \in \mathcal{E}$.
    
    Thus, all algebras containing $\calC$ will contain $\calA(\calC)$, and thus, $\calA(\calC)$ is the smallest algebra containing $\calC$.
    \end{proof}
    
    (ii)
    \begin{proof}
        \textbf{Part 1:} WTS: $\gensig{\calC} \subset \gensig{\calA(\calC)}$ \\
        $\calC \subset \calA(\calC) \: \implies \gensig{\calC} \subset \gensig{\calA(\calC)}$


        \textbf{Part 2:} WTS: $ \gensig{\calA(\calC)} \subset \gensig{\calC}$ \\
        $\calA(\calC) \subset \gensig{\calC} \implies \gensig{\calA(\calC)} \subset \gensig{\gensig{\calC}} = \gensig{\calC}$

        Thus, $\gensig{\calC} = \gensig{\calA(\calC)}$
    \end{proof}
    \end{solution}

%%%%%%%%%%% Question 6 %%%%%%%%%%%%%%%%%%%%    
    \question Chapter 1, Problem 1.20
    
    Let $\mu^*$ be as in (3.1) of Section 1.3: 
    
    Given a measure $\mu$ on a semialgebra $\calC \subset \power(\Omega)$, the outer measure induced by $\mu$ is the function $\mu^*$, defined on $\power(\Omega)$ as:
            \[\mu^*(A) = \inf\left\{\sum_{n = 1}^\infty \mu(A_n): \{A_n\}_{n \geq 1} \subset \calC, A \subset \bigcup_{n \geq 1} A_n\right\}\]
            
    Verify:
    \begin{itemize}
        \item[(3.4)] $\om (\emptyset) = 0$
        \item[(3.5)] $A \subset B \: \implies \om(A) \leq \om(B)$
        \item[(3.6)] For any $\{A_n\}_{n \geq 1} \subset \power(\Omega)$, 
        \[\om\left(\bigcup_{n \geq 1} A_n\right) \leq \sum_{n = 1}^\infty \om(A_n).\]
    \end{itemize}
    (\textbf{Hint}: Fix $0 <\epsilon<\infty$. If $\mu^*(A_n) < \infty$ for all $n \in \nat$, then find, for each $n$, a cover $\{A_{nj}\}_{j\geq 1} \subset \calC$, $A_n \subset \bigcup_{n\geq 1}A_{nj}$ such that $\sum^\infty_{j=1} \mu(A_{nj}) \leq \mu^*(A_n) + \frac{\epsilon}{2n}$.)

    \begin{solution}

    Let $\calC$ be a semialgebra and and $\mu$ be a measure defined on $\calC$. Define $\calA = \calA(\calC)$ as the smallest algebra generated by $\calC$. Let $\om$ be the outer measure induced by $\mu$

    (3.4)
        \begin{proof}
            \textbf{Part 1:} WTS $\emptyset \in \calC$

            Let $A \in \calC$ be given. Thus there exists $B_1, ... , B_k \in \calC$ for some $k < \infty$ such that $B_i \cap B_j = \emptyset$ for $i \neq j$ and $\cup_{i = 1}^k B_i = A^c$

            Since $\calC$ is closed under finite intersection for some $i \neq j$, $B_i \cap B_j = \emptyset \in \calC$. 

            Since $\mu$ is a measure, $\mu(A) \geq 0$ for all $A \in \calC$. Thus, $\om(A) \geq 0$ as well. 

            By definition, $\mu(\emptyset) = 0$. Define $A_n \equiv \emptyset$ for all $n \geq 1$. $\{A\}_{n \geq 1} \subset \calC$ and $\emptyset \subset \cup_{i = 1}^n A_n$. \\
            $\sum_{i = 1}^\infty A_n = 0$. 0 is the smallest value that $\om$ can measure, so $\om(\emptyset) = 0$.
        \end{proof}
    (3.5)
        \begin{proof}
        Suppose $A \subset B$.

        Thus, 
        \[\left\{\sum_{n = 1}^\infty \mu(A_n): \{A_n\}_{n \geq 1} \subset \calC, A \subset \bigcup_{n \geq 1} A_n\right\} \subset \left\{\sum_{n = 1}^\infty \mu(B_n): \{B_n\}_{n \geq 1} \subset \calC, B \subset \bigcup_{n \geq 1} B_n\right\}\]
        Thus, 
        \[\om(A) = \inf\left\{\sum_{n = 1}^\infty \mu(A_n): \{A_n\}_{n \geq 1} \subset \calC, A \subset \bigcup_{n \geq 1} A_n\right\}\]
        \[\subset \inf\left\{\sum_{n = 1}^\infty \mu(B_n): \{B_n\}_{n \geq 1} \subset \calC, B \leq \bigcup_{n \geq 1} B_n\right\} = \om(B)\]
        \end{proof}
\end{solution}
\newpage 
\begin{solution}
    (3.6)
        \begin{proof}
            Let $\{A_n\}_{n \geq 1} \subset \power(\Omega)$ be given. Let $0 < \epsilon < \infty$. \\
            Because, 
            \[\mu^*(A_n) = \inf\left\{\sum_{n = 1}^\infty \mu(A_n): \{A_{nj}\}_{n \geq 1} \subset \calC, A_n \subset \bigcup_{n \geq 1} A_{nj}\right\}\]
        For each $A_n \in \{A_n\}_{n \geq 1},$ there exist $\{A_{nj}\}_{j \geq 1} \subset \calC$ such that $A_n \subset \bigcup_{j \geq 1} A_{nj}$ and $\sum_{j \geq 1}\mu(A_{nj}) \leq \om(A_n) + \frac{\epsilon}{2^n}$.

        Thus 
        \begin{equation}
            \sum_{n \geq 1} \sum_{j \geq 1}\mu(A_{nj}) \leq \sum_{n \geq 1} \om(A_n) + \sum_{n \geq 1} \frac{\epsilon}{2^n} = \sum_{n \geq 1} \om(A_n) + \epsilon.
        \end{equation}

        Since $A_n \subset \bigcup_{j \geq 1} A_{nj}$ for all $n \geq 1$, $\bigcup_{n\geq 1} A_n \subset \bigcup_{n\geq 1}\bigcup_{j \geq 1} A_{nj}$

        By the definition of an outer set, $\om(\cup_{n \geq 1}A_n) \leq \sum_{n \geq 1}\sum_{j \geq 1} A_{nj}$. 

        By, (\theequation), $\om(\cup_{n \geq 1}A_n) \leq \sum_{n \geq 1}\sum_{j \geq 1} A_{nj} \leq  \sum_{n \geq 1} \om(A_n) + \epsilon$. 

        Since this is true for all $0 < \epsilon < \infty$, $\om(\cup_{n \geq 1}A_n) \leq \sum_{n \geq 1} \om(A_n).$
        \end{proof}
    \end{solution}

%%%%%%%%%%% Question 7 %%%%%%%%%%%%%%%%%%%%    
    \question Chapter 1, Problem 1.24 \\
    Verify that $\calC_2$, defined in (3.11), is a semialgebra. 
    \begin{solution}
    $\calC= \{(a, b]: -\infty < a, b < \infty\} \cup  \{(a, \infty): -\infty < a < \infty\} $    \\
    $\calC_2 = \{I_1 \times I_2: I_1, I_2 \in \calC\}$

    \begin{proof}
        \textbf{Part 1:} WTS: $A, B \in \calC_2 \implies A \cap B \in \calC_2$ \\
        Let $A, B \in \calC_2$ be given.\\ 
        Thus, there exist $I_{1, A}, I_{1, B}, I_{2, A}, I_{2, B} \in \calC$ such that $A = I_{1, A} \times I_{2, A}$ and $B = I_{1, B} \times I_{2, B}$. Since $\calC$ is a semi-algebra (given to us in the textbook), $I_{1, A} \cap I_{1, B} \in \calC$ and $I_{2, A} \cap I_{2, B} \in \calC$, which implies that $A \cap B \in \calC_2$.

        \textbf{Part 2:} WTS: For all $A \in \calC_2$, there exist sets $B_1, B_2, ..., B_k \in \calC_2$ for some $1 \leq k < \infty$ such that $B_i \cap B_j = \emptyset$ for $i \neq j$ and $A^c = \cup_{i =1}^k B_i$

        Let $A \in \calC_2$ be given. Thus, there exist $I_1, I_2 \in \calC$ such that $A = I_1 \times I_2$. \\
        Since $\calC$ is a semialgebra, there exist $B_{1, 1}, B_{1, 2}, B_{1, k_1}, B_{2, 1}, B_{2, 2}, B_{2, k_2} \in \calC$ such that $B_{1, i} \cap B_{1, j}, B_{2, i} \cap B_{2, j} = \emptyset$ for $i \neq j$, and $\bigcup_{i \leq k_1} B_{1, i} = I_1^c$ and $\bigcup_{i \leq k_2} B_{2, i} = I_2^c$ for some $0 < k_1, k_2 < \infty$. Without loss of generality, suppose $k_1 > k_2$. Thus, define $B_{2, i} \equiv \emptyset$ for all $i > k_2$.

        Define: $B_i \equiv B_{1, i} \times B_{2, i}$ for all $1 \leq i \leq k_1$. It is clear to see that $B_i \cap B_j = \emptyset$ for $i \neq j$, and $A^c = \cup_{i =1}^{k_1} B_i$.

        Thus, $\calC_2$ is a semialgebra
    \end{proof}
    \end{solution}

%%%%%%%%%%% Question 8 %%%%%%%%%%%%%%%%%%%%    
    \question Extra Problem 

    Consider the following function: \\
    $F(x) = x \ibold\{0 < x < 1 \} + \ibold\{x \geq 1\}, \: x \in \real$
    \begin{itemize}
        \item[(i)] Is this $F(\cdot)$ non-decreasing, right-continuous? Check.
        \item[(ii)] Define suitably a Lebesgue-Stieltjes measure $\mu_F$ on a suitable semialgebra $\calC$. Do not have to prove $\calC$ is a semi-algebra or $\mu_F$ is a measure.
        \item[(iii)] Prove that $\sigma \langle \calC\rangle =\calB(\real)$. (Can use the fact that $\calB(\real) = \sigma \langle \{\textrm{open sets}\}\rangle  = \sigma \langle \calO_1\rangle  =\sigma \langle \calO_2\rangle =\sigma \langle \calO_3\rangle  =\sigma \langle \calO_4\rangle $ as in 1.10 )
        \item[(iv)] Argue that there exist a $\mu$ on $(\real, \calB(\real))$ that extends $\mu_F$ defined above. Do not need to prove this step - just state a result for your answer.
        \item[(v)] Is this extension unique ? (i.e if $\mu'$ is another extension of $\mu_F$, then is it true that $\mu(A) = \mu'(A)$, for all $A\in \calB(\real)$ ?). Justify your answer (no proof needed). What is the “common name” for this measure $\mu$?
 
    \end{itemize}
    \begin{solution}
    (i) 
    \begin{proof}
        \textbf{Part 1}: $F(\cdot)$ is non-decreasing \\
        It is clear to see that this is true since for all $x_1 < x_2$, $F(x_1) \leq F(x_2)$

        \textbf{Part 2:} $F(\cdot)$ is right continuous \\
        Let $x_0 \in \real$ be given. \\
        \underline{Case 1:} $x_0 \leq 0$ \\
        $\lim\limits_{x \to x_0} F(x) = 0 = F(x_0).$\\
        Thus, $F(\cdot)$ is right continuous when $x \leq 0$. 

        \underline{Case 2:} $0 < x_0 < 1$ \\
        $\lim\limits_{x \to x_0} F(x) = x_0 = F(x_0)$. \\
        Thus, $F(\cdot)$ is right continuous when $0 < x < 1$

        \underline{Case 3:} $x_0 \geq 1$ \\
        $\lim\limits_{x \to x_0} F(x) = 1 = F(x_0)$. \\
        Thus, $F(\cdot)$ is right continuous when $ x \geq 1$

        Thus, $F(\cdot)$ is non-decreasing, right-continuous. 
    \end{proof}

    (ii) \\
    $\calC \equiv \{(a, b]: -\infty \leq a \leq b < \infty \} \cup \{(a, \infty): -\infty \leq a < \infty\}$ \\
    $\mu_F((a, b]) = F(b+) - F(a+) \\ \mu_F((a, \infty)) = F(\infty) - F(a+)$

    (iii) 
    \begin{proof}
    For all $b \in \real$, $(-\infty, b) = \bigcup_{n \geq 1} (-\infty, b + \tfrac{1}{n}] \: \implies \: \calO_2 \subset \calC \\ \implies \: \calB(\real) = \gensig{\calO_2} \subset \gensig{\calC}$

    For $\bigcup_{n \geq 1} (a, b + \tfrac{1}{n}) = (a, b] \: \implies \: \calC \subset \gensig{\calO_1} \: \implies \: \gensig{\calC} \subset \gensig{\gensig{\calO_1}} = \gensig{\calO_1} = \calB(\real)$

    Therefore, $\gensig{\calC} = \calB(\real)$
    \end{proof}

    (iv)

    We can use Caratheodory's Extension Theorem. $\mu_F$ is defined on $\calC$, so $\calC \subset \mathcal{M}_\om$ ($\om$ is the outer measure induced by $\mu_F$). Thus, $\gensig{\calC} = \calB(\real) \subset \mathcal{M}_\om$. Thus, the Lebesgue-Stieltjes measure extends to $\calB(\real)$.

    (v) Yes, by the Uniqueness theorem (1.3.6), since $\mu_F$ is a $\sigma$-finite measure on semialgebra $\calC$, per (iv), $\mu_F$ is unique. This is the uniform (0,1) measure on the real line.  

    \end{solution}

\end{questions}
\end{document}